\section{技术选型}

本章说明“鱼你有图”项目的技术栈、选型理由和工程适配关系。选型原则是采用成熟、稳定、便于部署和便于团队维护的技术组合，使地图服务、空间数据、推荐模型、AI Agent、社区内容和用户数据能够形成统一的工程闭环。

\subsection{技术栈总览}

\begin{table}[H]
  \centering
  \caption{技术栈总览}
  \label{tab:tech-stack}
  \begin{tabularx}{\textwidth}{p{3cm}|p{4.6cm}|X}
    \toprule
    \textbf{层次} & \textbf{技术选型} & \textbf{说明} \\
    \midrule
    前端框架 & Vue 3 + Vite & 构建移动端优先的单页应用，支撑地图、社区、记录、教程和个人中心等页面 \\
    \hline
    UI 组件库 & Element Plus + 图标组件 & 提供表单、弹窗、标签、按钮、图标等基础控件，保证页面风格统一 \\
    \hline
    地图能力 & 高德 JS API 2.0 + 高德 Web Service & 前端负责地图渲染，后端封装 POI、天气、地理编码等服务调用 \\
    \hline
    后端框架 & Python 3.11 + FastAPI + Pydantic 2 & 提供类型清晰的 RESTful API，适合规则推荐、AI Agent 和第三方服务封装 \\
    \hline
    数据库 & PostgreSQL + PostGIS & 保存业务数据和空间数据，支撑附近查询、空间索引、禁钓区域判断和推荐排序 \\
    \hline
    数据访问与迁移 & SQLAlchemy 2 + Alembic & 统一数据库模型、查询逻辑和表结构版本管理 \\
    \hline
    缓存与会话 & Redis & 缓存天气、热点 POI、推荐结果、登录会话和限流计数 \\
    \hline
    对象存储 & OSS/S3 兼容对象存储 & 存放鱼获图片、帖子图片、头像、教程视频封面和报告长图 \\
    \hline
    AI 服务 & DeepSeek API & 生成个性化出钓建议、推荐解释和复盘摘要，不直接决定排序 \\
    \hline
    网关与部署 & Nginx + HTTPS & 托管前端静态资源，反向代理 \texttt{/api}，处理 HTTPS 和访问日志 \\
    \hline
    测试与质量 & Pytest + 前端构建检查 & 覆盖后端服务、数据模型、推荐规则、权限控制和前端构建质量 \\
    \bottomrule
  \end{tabularx}
\end{table}

\subsection{前端选型理由}

前端采用 Vue 3 + Vite。Vue 组件模型清晰、学习成本适中、生态成熟，适合构建地图首页、钓点详情、社区信息流、发布表单、出钓记录、教程列表和个人报告等页面。Vite 提供快速构建和模块化开发能力，便于团队在页面交互、地图组件和 API 调用之间保持清晰边界。

Element Plus 用于基础控件和交互组件，能够减少重复开发，使页面风格、表单校验、弹窗交互和状态提示保持一致。地图类页面由高德 JS API 承接，普通业务页面由 Vue 组件承接，二者通过统一状态和 API 服务层协作。

\subsection{后端选型理由}

后端采用 FastAPI。FastAPI 与 Pydantic 结合后，可以直接定义 POI、记录、帖子、推荐、天气、用户偏好等模型，接口参数和响应结构清晰，适合投标文件、接口文档和代码实现保持一致。Python 生态对规则模型、统计分析、AI API 调用和数据处理友好，能够支撑推荐 Agent 的快速迭代。

后端分为 API 路由层、业务服务层、数据访问层和第三方服务封装层。API 路由层处理 HTTP 请求和参数校验；业务服务层承接推荐、记录、社区、报告和权限逻辑；数据访问层通过 SQLAlchemy 操作 PostgreSQL；第三方封装层统一接入高德、DeepSeek、对象存储和通知服务。

\subsection{数据库选型理由}

数据持久化采用 PostgreSQL + PostGIS。PostgreSQL 具备成熟的事务、索引、约束和备份能力，适合保存用户、记录、帖子、评论、活动、订单等关系明确的数据。PostGIS 提供空间类型、空间索引和空间函数，能够直接支撑附近钓点查询、禁钓区叠加判断、水域范围筛选和位置脱敏等 GIS 场景。

相比 MySQL，PostgreSQL + PostGIS 在空间数据处理方面更适合本项目；相比 MongoDB，关系型数据库更便于维护用户、记录、帖子、活动报名、互动数据之间的关联一致性。对于推荐结果、报告快照和第三方摘要等半结构化内容，PostgreSQL 的 \texttt{jsonb} 字段能够兼顾灵活性和查询能力。

\subsection{AI 与推荐选型理由}

推荐系统采用规则模型与 AI Agent 分工协作。规则模型负责可解释排序、合规拦截和风险降权，避免推荐结果不可验证；AI Agent 负责解释推荐原因、结合用户偏好生成出钓建议、形成复盘摘要和教程建议。该设计既能体现智能化服务能力，也能保证合规、安全和可审计。

DeepSeek API 用于自然语言生成。后端只向大模型传递结构化上下文和授权画像摘要，不传递完整隐私原文。大模型输出必须经过格式校验和风险约束，不允许直接覆盖规则排序结果。

\subsection{部署与运维选型理由}

Nginx 负责前端静态资源托管、HTTPS 终止、反向代理和访问日志。Redis 用于缓存热点数据、天气快照、推荐结果和限流计数，降低数据库和第三方服务压力。对象存储用于保存图片、视频、头像和报告长图，数据库仅保存文件元数据和访问地址。

Alembic 负责数据库迁移脚本管理，保证不同环境表结构一致；Pytest 用于后端单元测试和接口测试；前端构建检查用于保证页面资源可正常打包。通过这些技术组合，系统能够支撑课程交付、线上运行和生产部署。

